% !TeX root = ../thesis.tex
\chapter{Kurzfassung}

Im betrachteten Schulungsprozess bei der BMW Group wird die Teilnahme an internen und externen Weiterbildungen auf physischen Teilnehmerlisten dokumentiert, die anschließend gescannt oder fotografiert als \ac{pdf}-Dokumente archiviert und bislang manuell in das betriebliche Schulungssystem übertragen werden. Bei einem historischen Dokumentbestand von rund 6\,000 internen \acp{pdf} ist diese manuelle Nacherfassung zeitaufwendig und fehleranfällig. Zugleich stellen heterogene Tabellenlayouts, falsch gedrehte Scan-Seiten, mehrseitige Listen ohne wiederholte Spaltenüberschriften, gemischte Maschinen- und Handschrift sowie die strengen Datenschutzvorgaben von BMW und begrenzte lokale Grafikspeicherressourcen (\mbox{16\,GiB} \ac{vram}) hohe Anforderungen an eine automatisierte Verarbeitung.

Ziel dieser Bachelorarbeit ist die systematische Konzeption, Implementierung und empirische Bewertung \ac{ki}-basierter Extraktionsarchitekturen für diesen konkreten Anwendungsfall bei BMW zur strukturierten Gewinnung der fünf Zielfelder \emph{Vorname}, \emph{Nachname}, \emph{Personalnummer} (\emph{Benutzer} bzw.\ \emph{User~\ac{id}}), \emph{Abteilung} (\emph{Organisation/Kurzzeichen}) sowie des visuellen \emph{Vorhandenseins einer Unterschrift} (ausschließlich als Sichtbarkeitsprüfung ohne biometrische Verifikation). Hierzu wurde ein modulares, vollständig lokal ausführbares Evaluationsframework in Python entwickelt, das fünf ausgewählte Extraktionsverfahren (\texttt{start\_testframework\_5.cmd} / \texttt{configs/benchmark-five.yaml}) unter kontrollierten Bedingungen vergleicht: (1)~eine klassische, regel- und vorlagenbasierte \ac{ocr}-Baseline mit Tesseract ohne Sprachmodell (\texttt{ocr\_rules}), (2)~Tesseract-\acs{ocr} sowie (3)~PaddleOCR/PP-StructureV3 jeweils gekoppelt mit dem lokalen Text-\ac{llm} \texttt{Qwen3-4B-Instruct-2507} (\texttt{ocr\_llm} bzw.\ \texttt{paddle\_llm}), (4)~ein direkt auf dem Seitenbild arbeitendes, \acs{ocr}-freies \ac{vlm} auf Basis von \texttt{Qwen3-VL-4B-Instruct} (\texttt{vlm}) sowie (5)~ein kontrolliertes Hybridverfahren (\texttt{ocr\_vlm}), bei dem eine unabhängig vorgeschaltete PaddleOCR-Erkennung mit demselben \acs{vlm} kombiniert wird, sodass das Modell Seitenbild und \acs{ocr}-Zellbelege gemeinsam auswertet. Ergänzend wurde ein strikt isolierter \ac{api}-Evaluationspfad für den BMW-Firmen-PC entworfen.

Zur Sicherstellung der Extraktions- und Evaluationsgüte umfasst das Framework eine zweistufige automatische Seitenorientierung (\ac{osd} mit Vierlagen-\acs{ocr}-Fallback), eine zerstörungsfreie Hintergrundnormalisierung, ein verlustfreies Beobachtungsregister, eine geometrisch abgesicherte Spaltenvererbung für kopflose Folgeseiten, eine zeilenindividuelle Vor- und Nachnamen-Trennung, eine strikte Evidenz- und Positionsvalidierung sowie eine mehrdimensionale Metrik- und Fehleranalyse-Pipeline. Auf den bisher untersuchten realen Entwicklungsstichproben (ein zweiseitiges gemischtes Dokument mit 72 Referenzfeldern, ein Dreidokumentvergleich mit 184 Referenzfeldern sowie ein Zehndokumentlauf mit 620 Referenzfeldern) erreichen \texttt{paddle\_llm} (51,09\,\% exakter Feld-Recall, 60,65\,\% exakter Feld-$F_1$ im Dreidokumentlauf) und \texttt{ocr\_vlm} (47,83\,\% exakter Feld-Recall, nach abgesicherter Folgeseitenreparatur 52,17\,\%) die höchste Extraktionsqualität. Die kausale Fehleranalyse zeigt, dass ein erheblicher Teil der scheinbaren Extraktionsverluste nicht auf reiner Erkennungsblindheit beruht, sondern durch seitlich gedrehte Scans, kopflose Folgeseiten, Tokenbudget-Überschreitungen und Diskrepanzen zwischen externen Kursreferenzlisten und dem tatsächlich sichtbaren Dokumentinhalt entsteht. Eine abschließende, konfirmatorische Evaluation auf einer unabhängig annotierten Pilot- und Teststichprobe ist im Framework vorbereitet und bildet den unmittelbaren Folgeschritt.

\textbf{Schlagwörter:} Informationsextraktion, Teilnehmerlisten, Optical Character Recognition (\acs{ocr}), Large Language Models (\acs{llm}), Vision-Language-Modelle (\acs{vlm}), Document AI, Lokale \acs{ki}-Evaluation
